\section{Methodology}\label{sec:method}
The study is an exploratory development analysis. V7-B was shaped by seeing these data, no sealed or holdout split exists, and the confirmatory gate K5 (at least 60 projects, at least 250 independent divergent pairs, kappa of at least 0.7, clean leakage audit) is not met; under the pre-declared analysis plan the study is therefore exploratory. The plan (project-level inference, project-cluster bootstrap, sign-flip tests, Holm correction, a +10pp practical threshold, power recomputed from the actual design) was written before any held-out result existed, but it was applied here only to development data. Pipeline: run V7-B on 321 pairs through its command-line interface (spec configuration except per-run limits of 2~s CPU and 5~s wall, budget 60 CPU-s per pair, at most 12 inputs per pair, four workers); run the baselines on identical inputs; build one table of (pair, system) rows; analyse. The statistics in Section~\ref{sec:results} are from the post-fix rerun. Gates: K1 asks for at least 100 runnable real pairs from at least 8 projects; K2 for determinism and activation floors under proper isolation; K3--K5 are downstream and were not reached.
\dotfig{split_pipeline}{Development (current) versus planned DEV-N, sealing, holdout and confirmatory stages; the latter do not exist}{fig:split}
